% \documentclass[12pt]{article}
\documentclass[12pt]{book}
\hbadness=10000                       % Suppress Warnings (Horizontal)
\vbadness=10000                       % Suppress Warnings (Vertical)


% ------------------------------------------------------------------------------
% Required Packages for preamble/font.tex
% ------------------------------------------------------------------------------
\makeatletter
% Font encoding and families
\@ifpackageloaded{fontenc}{}{\usepackage[T1]{fontenc}}   % better font encoding
\@ifpackageloaded{lmodern}{}{\usepackage{lmodern}}       % Latin Modern — scalable replacement for CM

% Typography helpers
\@ifpackageloaded{lettrine}{}{\usepackage{lettrine}}      % dropped capitals
\@ifpackageloaded{lipsum}{}{\usepackage{lipsum}}          % dummy text (testing)
\makeatother

% ------------------------------------------------------------------------------
% Documentation and Usage Examples: lipsum and lettrine
% ------------------------------------------------------------------------------
% DESCRIPTION
%   The lipsum package provides dummy text for testing typographic layout.
%   The lettrine package produces dropped capitals at the beginning of a paragraph.
%
% FEATURES
%   - lipsum: generate one or more paragraphs for quick layout checks
%   - lettrine: control dropped capital lines, spacing, and scaling
%
% USAGE: lipsum
%   Basic:    \lipsum                  % generates a default paragraph
%   Paragraph: \lipsum[1]              % first paragraph
%   Range:     \lipsum[2-3]            % paragraphs 2 to 3
%   Words:     \lipsum[1][1-5]         % words 1–5 from paragraph 1
%
% USAGE: lettrine
%   Basic:     \lettrine{T}{his} starts a paragraph with a dropped “T”
%   Lines:     \lettrine[lines=3]{T}{his}    % drop across 3 lines
%   Spacing:   \lettrine[findent=2pt,nindent=0pt]{T}{his}
%   Scaling:   \lettrine[lines=3,loversize=0.15]{T}{his}
%   Slope:     \lettrine[hang=1]{T}{his}     % tighter hang of text under the drop cap
%
% EXAMPLES (copy into your document body to preview)
%   % Example 1: Simple lipsum
%   % \lipsum[1]
%
%   % Example 2: Variable Length lipsum
%   % \lipsum[2-3]
%
%   % Example 3: Drop cap with default settings
%   % \lettrine{L}{orem} ipsum dolor sit amet, consectetur adipiscing elit.\lipsum[2]
%
%   % Example 4: Three-line drop cap with adjusted size
%   % \lettrine[lines=3,loversize=0.2]{D}{olor} sit amet, consectetur adipiscing elit.\lipsum[3]
%
%   % Example 5: Fine-tuned indents for tighter text wrap
%   % \lettrine[lines=3,findent=1.5pt,nindent=0pt]{T}{ypography} testing block.\lipsum[4]
%
% NOTES
%   - Use \par after a lettrine paragraph if spacing looks off in complex layouts.
%   - Pair lettrine with a serif body font for best aesthetics.


\input{../preamble/font-korean.tex}
\input{../preamble/layout-letter.tex}
% ------------------------------------------------------------------------------
% Required Packages for preamble/reference-bibliography.tex
% ------------------------------------------------------------------------------
% Load packages conditionally
\makeatletter
\@ifpackageloaded{hyperref}{}{\usepackage{hyperref}} % helpful for TOC/citations/links
\makeatother

% Bibliography: biblatex with biber backend (works with pdflatex + biber)
\usepackage[
  backend=biber,           % use biber for processing
  style=numeric,           % bibliography style
  citestyle=authoryear,    % in-text citations
  sorting=ynt,             % year-name_of_auther–title order
  % maxcitenames=2,          % limit authors in citations
  % bibencoding=utf8         % if your .bib is UTF-8
]{biblatex}
\addbibresource{references.bib}

% ------------------------------------------------------------------------------
% USAGE
% ------------------------------------------------------------------------------
% Add citations in your document:
%   % \cite{key}   \parencite{key}   \textcite{key}
%   % Multiple: \parencites{key1}{key2}
%
% Print the bibliography (choose one of the styles):
%   % \clearpage
%   % \printbibliography[
%   %   heading=bibintoc,   % unnumbered entry in TOC
%   % %   title={References}
%   % ]
%
%   % \clearpage
%   % \printbibliography[
%   %   heading=bibnumbered,   % numbered entry in TOC
%   % %   title={References}
%   % ]
\input{../preamble/reference-index.tex}

\usepackage[indent]{parskip}

% page# location change to bottom like << articles >>
\usepackage{fancyhdr}
\pagestyle{fancy}
\fancyhf{}                        % clear all
\renewcommand{\headrulewidth}{0pt}
\renewcommand{\footrulewidth}{0pt}
\fancyhfoffset[L]{0pt}            % remove left offset
\fancyhfoffset[R]{0pt}            % remove right offset
\cfoot{\thepage}                  % page number centered



\usepackage{graphicx}

\newcommand{\VF}{\vspace*{\fill}}
\newcommand{\HF}{\hspace*{\fill}}
\newcommand{\tab}{\hspace*{2em}}



























\begin{document}
  %=======================================%
  %== Front Matter (roman page numbers) ==%
  %=======================================%
  \frontmatter
  \let\oldcleardoublepage\cleardoublepage
  \let\cleardoublepage\clearpage % redefines frontmatter's \cleardoublepage to a simple \clearpage
  
  
  \onesidedpaper
  % The \vspace*{\fill} command adds flexible vertical space. 
% Placing it before and after the text will center the content vertically on the page.
\vspace*{\fill}

\begin{center}
  % \noindent prevents any indentation at the beginning of the line.
  \noindent
  % The desired text for the blank page.
  \textit{This page intentionally left blank.}
\end{center}

% Adds flexible vertical space to push the content to the vertical center.
\vspace*{\fill}
  % Title page  
%================ Metadata ======================
\newcommand{\BookTitle}{Math Made Clear}
\newcommand{\BookSubtitle}{Optional Subtitle}
\newcommand{\BookAuthor}{Jueon Kim, Yeoeun Kim}
\newcommand{\BookEditor}{Jaehoon Song}
\newcommand{\BookInstitution}{Greater Youth Collaborative Opus (GYCO)}
\newcommand{\PublicationYear}{\the\year}
%================ Assets ======================
\newcommand{\BookEmblem}{./emblems/emblem.png} % Path to emblem image file
\begin{titlepage}
  \centering
  % Top fill for vertical centering
  \vspace*{\fill}
  % Main title block
  {\bfseries\Huge\BookTitle\par}
  \vspace{0.5em}
  % {\Large\BookSubtitle\par}
  \vfill
  \vspace{2em}
  % Author and editor
  {\Large Author: \BookAuthor\par}
  % {\Large Editor: \BookEditor\par}
  \vspace{2em}
  % Optional logo at bottom (uncomment to use)
  \includegraphics[width=0.65\textwidth]{\BookEmblem}
  {\normalsize \par}
  \vspace{2em}
  % Publisher and year
  {\normalsize \copyright\ \PublicationYear\ \BookInstitution\par}
  {\normalsize All rights reserved.\par}
  % Bottom fill for vertical centering
  \vspace*{\fill}
\end{titlepage}
  \newpage
% endorsements / praise (if any)
\section*{Endorsements}
\subsection*{추천의 글}
이 책은 우리 한글학교 학생들의 배움을 향한 사랑과 정성이 모여 탄생한 소중한 열매입니다. 
특히, 자원봉사 학생이 시간을 내어 영어 수학 용어를 한국어로 정리한 이 책은 또래 학생들을 
돕고자 하는 따뜻한 마음에서 시작되었습니다. 이 책을 통해 학생들이 수학을 조금 더 친근하게 
느끼고 나아가 한국어로 학문을 접하는 기쁨을 누리기를 바랍니다. 작은 노력이 큰 빛이 되어 
다른 학생들의 길을 밝혀주듯이 이 책이 많은 학생들에게 든든한 길잡이가 되기를 기대합니다. 
정성과 사랑으로 이 귀한 작업을 완성한 김주언, 김여은 학생에게 깊은 감사와 축복의 마음을 
전합니다.

\HF\ \textbf{쟌스크릭한국학교 교장 황영희}
\newline

자신의 어려웠던 경험을 생각하며 같은 어려움을 겪고 있는 학생들을 돕기 위해 수학용어 
책을 만들 생각을 하고, 생각에 그친 것이 아니라 바쁜 시간 속에서 책을 쓰기 위해 자료를 
수집하고 노력해 완성해 낸 너무 대견하고 자랑스러운 학생들입니다.

\inlinequote{배움과 나눔} 이라는 가치를 배워서 이타적인 삶을 어린 나이에 실천하고 있는 김주언, 
김여은 학생이 너무 자랑스럽습니다. 앞으로 본인들이 전공할 학문에서 더 깊이 있게 공부하고 
더 많은 사람들에게 지식과 마음을 나눌 수 있는 귀한 학생들이 될 것을 믿고 그 찬란하게 
빛날 앞길을 진심으로 격려하고 축복합니다.

\HF\ \textbf{비영리단체 GYCO 회장 김현정}
\newline
  \newpage
% Copyright page (on verso of title) MUST be LEFT page
\section*{Copyright}
\copyright\ 2025 Greater Youth Collaborative Orchestra (GYCO).

All rights reserved. No part of this book may be reproduced or used in any form 
without prior written permission of GYCO.

\subsection*{저작권}
\copyright\ 2025 Greater Youth Collaborative Orchestra (GYCO).

모든 권리 보유 이 책의 어떤 부분도 GYCO의 사전 허가 없이 복제하거나 사용할 수 없습니다.

  \tableofcontents% front 05
  % \listoffigures% front 06
  % \listoftables% front 07
  % Preface (by the author)
\newpage
\section*{Preface}
\addcontentsline{toc}{chapter}{Preface}
We wrote this book because we once struggled too. If you ever feel stuck, remember that you are not alone. Step by step, you'll grow stronger in both math and English. One day, you'll be the one helping others.



\subsection*{개인적인 메시지}
저희도 같은 어려움을 겪었기에 이 책을 만들었습니다. 힘들 때가 있어도 혼자가 아님을 기억하세요. 차근차근 연습하다 보면 수학과 영어 모두에서 성장할 것입니다. 언젠가는 여러분이 다른 사람을 도와줄 차례가 올 것입니다.

\blockquote{Thank you for learning with us. Don't give up and keep going-you've got this!} 

끝까지 함께해 주셔서 감사합니다. 포기하지 말고 계속 앞으로 나아가세요—여러분은 할 수 있습니다!
  
  % Acknowledgments (if not placed at the back)
  \section*{Acknowledgment}
  \addcontentsline{toc}{chapter}{Acknowledgment}
  We are deeply grateful to our loving parents, teachers, and friends who encouraged 
  and supported us throughout the process of creating this book. This project began 
  from our experience of struggling with different terms and units when we first 
  came to the United States, and we sincerely hope it will help Korean students 
  facing similar challenges.
  
  Special thanks to our parents and to GYCO for teaching us that learning and 
  sharing are the greatest values and for giving us chances to put them into practice.
  
  \subsection*{감사의 글}
  이 책을 만드는 과정에서 격려와 지지를 보내주신 사랑하는 부모님, 선생님들, 친구들에게 깊은 
  감사를 드립니다. 처음 미국에 왔을 때 용어와 단위가 한국과 달라 어려웠던 저희들의 경험에서 
  시작된 이 책이 같은 어려움을 겪고 있는 한국 친구들에게 도움이 되길 진심으로 바랍니다.
  
  늘 배움과 나눔을 가장 큰 가치로 가르쳐 주시고 실천하게 해주신 부모님과 GYCO에 특별한 
  감사를 전합니다.
  
  % Introduction
\newpage\section*{Introduction}
\addcontentsline{toc}{chapter}{Introduction}
This book was created to support students in Korean language schools and 
ESOL learners from Korea who face challenges with English. It is designed to 
strengthen both Korean language skills and a solid foundation in mathematics.

Each section introduces a core topic with clear explanations, followed by 
examples and practice problems. The main units include Fractions and Decimals, 
Factors and Multiples, and Ratios and Proportions.

Within each unit, you will find step-by-step guides, key math vocabulary used 
in the U.S., and real-world applications. To make the most of this book, read 
each section carefully, try solving the examples yourself, and learn the new 
vocabulary. Don't hesitate to revisit earlier sections if you need a refresher.

Learning math in a new language takes time and practice, but with patience 
and persistence, you can succeed. We hope this book becomes a valuable 
companion on your journey to mastering both math and Korean.

\subsection*{서문}
이 책은 한국어를 배우는 한글학교 학생들과 한국에서 미국으로 이주해 영어의 문제로 
어려움을 겪고 있는 ESOL 학생들을 돕기 위해 만들어졌습니다. 이 책은 여러분의 한국어 
실력을 향상시키는 동시에 수학에 대한 탄탄한 기초를 쌓을 수 있도록 구성되어 있습니다.

각 섹션은 핵심 주제를 다루며, 알기 쉽고 정확한 설명으로 시작해 예시와 연습 문제로 
이어집니다. 분수와 소수, 약수와 배수, 비와 비례식 등의 주요 단원으로 나뉩니다. 
각 단원에서는 단계별 가이드, 미국 수학 어휘 목록, 실제 응용 사례를 찾을 수 있습니다.

이 책을 최대한 활용하려면 각 섹션을 주의 깊게 읽고, 예시를 직접 풀어보며, 새로운 
어휘를 익히는 것이 좋습니다. 복습이 필요할 때는 이전 섹션을 다시 찾아보는 것을 주저하지 
마세요. 새로운 언어로 수학을 배우는 것은 시간과 연습이 필요합니다. 자기 자신에게 
인내심을 갖고 꾸준히 연습하면 충분히 잘할 수 있습니다. 우리는 이 책이 여러분의 
수학과 한국어를 함께 마스터하는 여정에 소중한 동반자가 되기를 희망합니다.
  









  
  %=======================================%
  %== Main Matter (arabic page numbers) ==%
  %=======================================%
  \mainmatter
  \let\cleardoublepage\oldcleardoublepage % revert redefines









  \twosidedpaper % Switch to the two-sided layout


  
  \chapter{Simulated One-Sided Layout}
  \lettrine{L}{orem ipsum} dolor sit amet, consectetur adipiscing elit. 
  This section visually appears as a one-sided document because the inner and 
  outer margins are equal on every page. s
  \lipsum[1-2]

  \lettrine{L}{orem ipsum} dolor sit amet \lipsum[1-2]



  \section*{Example - Line on this page only}
  \sidenote[2pt]{Note in margin dsa 
  dsad dsa dsa dsadsa dsad sadsadsadjkbsak dhsakjldskladb s sadsad.}

  \lettrine[lines=3]{L}{orem ipsum} dolor sit amet, consectetur adipiscing elit. 
  Praesent vel risus in nunc gravida venenatis. Curabitur euismod.
  \lipsum[1-2]


  Here is some normal text with a margin note.\sidenote{여기에는 한글 설명이 들어가야하는 자리입니다.}

  \section{Inside a float}

  \begin{figure}[ht]
    \centering
    \rule{3cm}{2cm}
    \caption{A box figure}
    Here is a figure caption. \sidenote{Margin note inside float.}
  \end{figure}

  \section{Inside a footnote}

  Some text with a footnote\footnote{여기 또한 마찬가지, 한글이 들어가야하는 자리입니다. 영어 단어나 문장 해석을 목적으로 제공되는 공간입니다.
  \lipsum[1] sad sad \sidenote[-10pt]{Margin note inside footnote.}} continuing the paragraph.






  \newpage
  \section*{Example - Line on this page only}

  \lipsum[1-2]
  \begin{table}[h!]
    \centering
    \caption{Sample Table of Student Scores} % Appears in List of Tables
    \begin{tabular}{|c|c|c|}
      \hline
      Student & Subject & Score \\
      \hline
      Alice   & Math    & 95    \\
      Bob     & Physics & 88    \\
      Carol   & English & 92    \\
      \hline
    \end{tabular}
    \label{tab:scores} % useful for referencing
  \end{table}

  \section{Inside a footnote}

  Some text with a footnote\footnote{Footnote text. 
  dsadsadsad 
  sadsad sadsad sadsad sad sad sadsad \lipsum[1-2] sad sad \sidenote{Margin note inside footnote.}} continuing the paragraph.



  Here is some normal text with a margin note.\sidenote{Note in margin dsa 
  dsad dsa dsa dsadsa dsad sadsadsadjkbsak dhsakjldskladb s sadsad.}






  \newpage
  \section*{Example - Line on this page only}
  \lipsum[1-2]



  \newpage
  \section*{Example - Line on this page only}
  \lipsum[1-2]






  \onesidedpaper % Start with the one-sided layout
  \section*{Part 1: Simulated One-Sided Layout}
  This section visually appears as a one-sided document because the inner and 
  outer margins are equal on every page. s
  \lipsum[1-2]

  \newpage
  \section*{Part 2: True Two-Sided Layout}
  The page layout has now switched to a true two-sided style. The margins on even and odd pages will be different.
  \lipsum[3-4]





    
  \section{Number theory example}
  Let $q$ be any prime divisor of \(N\). We make two observations:
  \index{prime divisor}
  \index{$q$@\texttt{$q$}} % put q in index (show as verbatim-like text)

  This is English text using the Latin font set by pdfLaTeX.
  \index{English text}

  \section{Korean examples}
  안녕하세요. 이것은 \texttt{kotex}로 컴파일된 한글입니다.
  \index{kotex}

  \subsection{Nanum family examples}
  안녕하세요! 이것은 나눔명조체입니다.
  \index{나눔명조}

  \textit{이것은 나눔손글씨입니다.}
  \index{나눔손글씨}

  \textbf{이것은 나눔바른고딕입니다.}
  \index{나눔바른고딕}

  \textbf{\textit{이것은 나눔손글씨입니다.}}
  \index{나눔손글씨!강조}

  \section{More sample entries}
  You can add hierarchical (sub) entries, e.g., Fonts!Nanum:
  \index{Fonts!Nanum Myeongjo}
  \index{Fonts!Nanum Gothic}


  Cite an authoritative book on algorithms: \textcite{noauthor_help_2010} shows a number of classic results.
  Also see \parencite{kim_ssat-sat_2003} for document preparation practices.

  











  %=======================================%
  %== Back Matter ==%
  %=======================================%  
  \backmatter
  % Appendix 
  \chapter*{Appendix}
  \addcontentsline{toc}{chapter}{Appendix}
  Jueon Kim and Yeoeun Kim are students who have l


  \clearpage
  \printbibliography[
    heading=bibintoc,   % unnumbered entry in TOC
    title={References}
  ]
  \clearpage
  \printindex                  % unnumbered, title from imakeidx


  % Introduction
  \chapter*{About the Author}
  \addcontentsline{toc}{chapter}{About the Author}
  Jueon Kim and Yeoeun Kim are students who have lived in Korea, China, and the 
  United States, where they experienced firsthand how language and culture shape 
  learning. Adapting to new environments and study methods brought many challenges, 
  but guided by their parents' teaching to \inlinequote{learn well and share generously,} 
  they embraced service as a core value.
  
  Through their nonprofit organization GYCO, they bring music to seniors, the 
  homeless, and cancer patients each month, and send books to refugees and students 
  in Africa to support education. After performances, they talk with many seniors 
  to help ease loneliness.
  
  Their commitment to service has inspired many peers, and this math book reflects 
  their own journey as Korean immigrants adapting to new languages and education 
  systems. They hope it helps students overcome language barriers, see math as 
  approachable, and use it as a tool for confidence and opportunity.
  
  \section*{저자 소개}
  김주언, 김여은 두 학생은 한국, 중국, 미국에서 생활하며 언어와 문화가 학습에 어떤 
  영향을 주는지 직접 경험한 학생입니다. 이 학생들은 변화된 환경에 적응해야 했고 새로운 
  문화와 학습 방법을 받아들이는 과정 속에서 어려움도 많았습니다.
  
  부모님께 배운 \inlinequote{잘 배워 잘 나누는 사람이 되라}는 가르침은 이들의 삶의 중요한 
  가치가 되었습니다. 비영리단체 GYCO에서 사회의 소외계층에게 사랑과 희망을 전하는 다양한 
  활동을 현재까지 꾸준하게 하고 있습니다.
  
  음악 연주를 직접 보러 올 수 없으신 사회의 연약한 계층(시니어, 노숙자, 암환자)에게 한 
  달에 한 번씩 연주로 찾아가고, 교육이 필요한 난민과 아프리카의 학생들에게 책을 보내 
  교육받을 수 있게 돕고 있습니다. 많은 시니어들에게 찾아가 연주가 끝나고 함께 이야기를 
  나누며 그들의 외로움을 덜어드렸고, 노숙자분들께 필요한 물품을 준비해서 나눠드리는 활동도 
  지속적으로 해오고 있습니다. 함께 봉사하는 학생들에게도 이들의 성실함과 열정과 나눔은 늘 
  도전이 되고 있습니다.
  
  이 수학책은 한국 이민자로서 새로운 언어와 교육 제도에 적응하면서 느낀 두 학생의 어려움을 
  바탕으로, 같은 길을 걷는 학생들이 수학을 더 쉽고 친근하게 느끼도록 돕기 위해 시작되었습니다.
  
  언어의 장벽으로 힘들어 하는 학생들이 수학을 어려움으로 느끼지 않고, 변화된 상황에서 뛰어난 
  수학 실력이 그들에게 자신감을 갖고 기회를 넓혀주는 도구가 되기를 희망하며 책을 만든 
  귀한 학생들입니다.




\end{document}